\chapter*{Conclusioni}
\addcontentsline{toc}{chapter}{Conclusioni}

In questo lavoro di tesi è stato affrontato il problema della protezione della privacy dei dati in contesti di pubblicazione di statistiche aggregate partendo dallo studio del framework della privacy differenziale proponendo la realizzazione di uno strumento software e discutendo una valutazione sperimentale.

Metodi come k-anonimity e l-diversity, pur offrendo protezione ragionevole quando considerati in modo isolato, si rivelano vulnerabili nel momento in cui un attaccante ha accesso a dati provenienti da fonti ausiliare, come social network, data breach o altre fonti pubbliche. La privacy differenziale nasce proprio per rispondere a questo tipo di minaccia e il lavoro svolto in questa tesi si è concentrato sul renderla concretamente utilizzabile.

Lo studio dei meccanismi alla base della DP, come calibrazione del rumore tramite sensitività, significato operativo dei parametri $\varepsilon$ e $\delta$, le proprietà di composizione del budget, ha rappresentato il fondamento necessario per passare dalla teoria all'aimplementazione. Senza un comprensione solida di questi aspetti sarebbe stato difficile progettare uno strumento che offrisse garanzie reali e non solo nominali.

\texttt{GoDP} è il risultato di questo percorso, si tratta di un'applicazione a linea di comando costruita su Privacy on Beam di Google, che permette di creare distribuzioni di dati differenzialmente private a partire da dataset in formato CSV e un file di configurazione yaml che descrive una serie di operazioni DP e scritto in una specifica yaml a corredo chiamata \texttt{DPYaml}. L'idea di base è stata quella di abbassare la barriera d'ingresso: un ricercatore o un piccolo team che ha bisogno di pubblicare statistiche aggregate su dati sensibili può evitare di implementare da zero meccanismi DP o a districarsi tra le API di un framework complesso. Con \texttt{GoDP}, è sufficiente descrivere le operazioni desiderate in un file YAML, specificare i parametri di privacy e lanciare il comando. Detto questo, l'architettura non è rigida: chi ha esigenze più specifiche può definire operazione personalizzate interagendo direttamente con i moduli interni dell'applicazione.

Uno degli aspetti che ha richiesto più attenzione in fase di sviluppo è stata la gestione del budget di privacy, GoDP fornisce un metodo semplice per destinare porzioni di budget maggiori o minori a una certa operazione a seconda della sua importanza, così come la possibilità di destinare più o meno budget alla selezione delle partizioni e all'aggregazione dei dati tramite la logica implementata nel modulo \texttt{budget} che applica queste suddivisioni garantendo l'assenza di errori che possono compromettere le garanzie offerte dalla DP.

L'esperimento condotto nel capitolo finale mette alla prova lo strumento in uno scenario concreto, un differencing attack su un dataset di dipendenti aziendali, dove un attaccante tenta di risalire allo stipendio di una vittima combinando i risultati di più query aggregate. Senza protezione DP l'attacco va a buon fine: la differenza di tre query aggregate è sufficiente per identificare il record della vittima; con GoDP a protezione del dataset, i risultati cambiano in modo significativo. Per $\varepsilon = 0.5$ il meccanismo di selezione delle partizioni rimuove completamente i risultati di due delle tre query, impedendo l'attacco ma allo stesso tempo azzerando l'utilità dei dati. Per $\varepsilon = 5$ l'attaccante ottiene un risultato con errore di circa il 50\% che rende l'inferenza sostanzialmente inutilizzabile. Per valori più permissivi come $\varepsilon = 10$ e $\varepsilon = 15$ l'errore si riduce, ma rimane rilevante a causa della composizione del rumore attraverso le tre query.

Questo risultato evidenzia come la scelta del valore di $\varepsilon$, e quindi del budget di privacy, dipenda largamente dal dominio dei dati e dal motivo per cui si vogliono rendere disponibili. In ambito medico, ad esempio, dove è imperativo evitare l'esposizione di un singolo dato, valori bassi di $\varepsilon$ sono necessari anche al costo di sacrificare precisione; per analisi interne meno sensibili valori più alti possono offrire un compromesso accettabile. La privacy differenziale non elimina questo problema decisionale che rimane tuttora ampiamente discusso ma lo rende esplicito e quantificabile\cite{epsilonExamples}\cite{howMuchIsEnough}.

Nella sua forma attuale \texttt{GoDP} offre caratteristiche interessanti che mirano a semplificare l'approccio alla differential privacy da parte di utenti non esperti, offrendo una gamma limitata di possibili operazioni che può essere estesa permettendo agli utenti di definire operazioni più complesse. La scelta di utilizzare Privacy on Beam (Apache Beam backend) come framework per implementare i meccanismi DP permette di portare il carico computazionale su runner distribuiti, il che renderebbe possibile elaborare dataset in modo altamente scalabile. Un altro miglioramento futuro possibile consiste nel rendere questo strumento adatto ad applicazioni del modello DP decentralizzato, portando l'applicazione dei meccanismi DP direttamente alla fonte della raccolta dati, eliminando la necessità di un curatore fidato.